\subsection{李代数相对于自同构的分解}

命题 12. 设 \(\mathfrak{g}\) 是一个李代数，a 是 \(\mathfrak{g}\) 的一个自同构。

\begin{enumerate}
\def\labelenumi{(\roman{enumi})}
\item
  对 \(\lambda, \mu \in k\) ，\([\mathfrak{g}^{\lambda}(a), \mathfrak{g}^{\mu}(a)] \subset \mathfrak{g}^{\lambda \mu}(a)\) ；特别地，\(\mathfrak{g}^1(a)\) 是 \(\mathfrak{g}\) 的一个子代数。
\item
  若 B 是 g 上在 a 下不变的对称双线性型，则 \(\mathfrak{g}^{\lambda}(a)\) 与 \(\mathfrak{g}^{\mu}(a)\) 关于 B 正交，如果 \(\lambda\mu \neq 1\) 。假定 B 非退化。那么，若 \(\lambda \neq 0\) ，则 B 在 \(\mathfrak{g}^{\lambda}(a) \times \mathfrak{g}^{1/\lambda}(a)\) 上的限制非退化。
\end{enumerate}

断言 (i) 与 (ii) 的前一半由把命题 2 (iii) 应用于合成律 \(\mathfrak{g} \times \mathfrak{g} \to \mathfrak{g}\) 及双线性型 B 得出。为证 (ii) 的后一半，可假定 \(k\) 是代数闭的。于是 \(\mathfrak{g} = \bigoplus_{\nu \in k} \mathfrak{g}^{\nu}(a)\) 。由前所述，\(\mathfrak{g}^{\lambda}(a)\) 与 \(\mathfrak{g}^{\nu}(a)\) 正交，如果 \(\lambda \nu \neq 1\) ；由于 B 非退化，可知 B 在 \(\mathfrak{g}^{\lambda}(a) \times \mathfrak{g}^{1/\lambda}(a)\) 上的限制也非退化。

推论. 假定 \(k\) 的特征为零且 \(\mathfrak{g}\) 是半单的。那么子代数 \(\mathfrak{g}^1(a)\) 满足引理 2 的条件 (1) 与 (2)；它在 \(\mathfrak{g}\) 中是约化的。

条件 (1) 由命题 12 的 (ii) 得出；条件 (2) 由第 1 节的命题 4 得出。

